% Source: https://artofproblemsolving.com/community/c3h2346064_msm_is_open_discuss_the_2020_amc_8_here
% Stable identity: aops-topic-2346064-q20
Tree 2 is $11$ m. Tree 1 is twice or half of it; $\dfrac{11}2$ is not an integer, so tree 1 is $22$ m. Likewise tree 3 is $22$ m. The average ends in $.2$, so the total $T$ of the five heights satisfies $\dfrac T5=k+0.2$ for an integer $k$, i.e. $T$ leaves remainder $1$ when divided by $5$.

Tree 4 is $11$ or $44$:
\begin{itemize}
\item tree 4 $=11$: then tree 5 is $22$ ($\dfrac{11}{2}$ is not an integer), $T=22+11+22+11+22=88$, average $17.6$;
\item tree 4 $=44$, tree 5 $=22$: $T=22+11+22+44+22=121$, average $24.2$;
\item tree 4 $=44$, tree 5 $=88$: $T=187$, average $37.4$.
\end{itemize}
Only the second case has an average ending in $.2$, so the average height is $24.2$ meters. The answer is $\boxed{\textbf{(B) }24.2}$.
